
\documentclass{article} % For LaTeX2e
\usepackage{iclr2026_conference,times}

% Optional math commands from https://github.com/goodfeli/dlbook_notation.
\input{math_commands.tex}

\usepackage{hyperref}
\usepackage{url}
\usepackage{algorithm}
\usepackage{algpseudocode}

%%%%%%%%%%%%%%%SAHANA ADDED THIS %%%%%%%%%%%
\usepackage{amsthm}
\usepackage{cleveref}
\newtheorem{theorem}{Theorem}[section]
\newtheorem{lemma}[theorem]{Lemma}
\newtheorem{claim}[theorem]{Claim}
\newtheorem{proposition}[theorem]{Proposition}
\newtheorem{corollary}[theorem]{Corollary}
\newtheorem{fact}[theorem]{Fact}
\newtheorem{example}[theorem]{Example}
\newtheorem{notation}[theorem]{Notation}
\newtheorem{observation}[theorem]{Observation}
\newtheorem{conjecture}[theorem]{Conjecture}

\newtheorem{innercustomgeneric}{\customgenericname}
\providecommand{\customgenericname}{}
\newcommand{\newcustomtheorem}[2]{%
  \newenvironment{#1}[1]
  {%
   \renewcommand\customgenericname{#2}%
   \renewcommand\theinnercustomgeneric{##1}%
   \innercustomgeneric
  }
  {\endinnercustomgeneric}
}

\newcustomtheorem{customtheorem}{Theorem}
\newcustomtheorem{customlemma}{Lemma}
\newcustomtheorem{customclaim}{Claim}
\newcustomtheorem{customproposition}{Proposition}
\newcustomtheorem{customcorollary}{Corollary}
\newcustomtheorem{customfact}{Fact}
\newcustomtheorem{customexample}{Example}
\newcustomtheorem{customnotation}{Notation}
\newcustomtheorem{customobservation}{Observation}
\newcustomtheorem{customconjecture}{Conjecture}
%%%%%%%%%%%%%%%%%%%%%%%%%%%%%%%%%%%%%%%%


\title{Greedy Hybrid PDE Solver with \\ Neural Operators and Traditional solvers}

% Authors must not appear in the submitted version. They should be hidden
% as long as the \iclrfinalcopy macro remains commented out below.
% Non-anonymous submissions will be rejected without review.

\author{Sahana Rayan \\
Department of Statistics\\
University of Michigan\\
Ann Arbor, MI 48104, USA \\
\texttt{\{srayan\}@umich.edu} \\
\And
Yash Patel \\
Department of Statistics\\
University of Michigan\\
Ann Arbor, MI 48104, USA \\
\texttt{\{yppatel\}@umich.edu} \\
\And
Ambuj Tewari \\
Department of Statistics\\
University of Michigan\\
Ann Arbor, MI 48104, USA \\
\texttt{\{tewaria\}@umich.edu} \\
}
% \author{Antiquus S.~Hippocampus, Natalia Cerebro \& Amelie P. Amygdale \thanks{ Use footnote for providing further information
% about author (webpage, alternative address)---\emph{not} for acknowledging
% funding agencies.  Funding acknowledgements go at the end of the paper.} \\
% Department of Computer Science\\
% Cranberry-Lemon University\\
% Pittsburgh, PA 15213, USA \\
% \texttt{\{hippo,brain,jen\}@cs.cranberry-lemon.edu} \\
% \And
% Ji Q. Ren \& Yevgeny LeNet \\
% Department of Computational Neuroscience \\
% University of the Witwatersrand \\
% Joburg, South Africa \\
% \texttt{\{robot,net\}@wits.ac.za} \\
% \AND
% Coauthor \\
% Affiliation \\
% Address \\
% \texttt{email}
% }

% The \author macro works with any number of authors. There are two commands
% used to separate the names and addresses of multiple authors: \And and \AND.
%
% Using \And between authors leaves it to \LaTeX{} to determine where to break
% the lines. Using \AND forces a linebreak at that point. So, if \LaTeX{}
% puts 3 of 4 authors names on the first line, and the last on the second
% line, try using \AND instead of \And before the third author name.

\newcommand{\fix}{\marginpar{FIX}}
\newcommand{\new}{\marginpar{NEW}}

% \iclrfinalcopy % Uncomment for camera-ready version, but NOT for submission.
\begin{document}


\maketitle

\begin{abstract}
The abstract paragraph should be indented 1/2~inch (3~picas) on both left and
right-hand margins. Use 10~point type, with a vertical spacing of 11~points.
The word \textsc{Abstract} must be centered, in small caps, and in point size 12. Two
line spaces precede the abstract. The abstract must be limited to one
paragraph.
\end{abstract}

\section{Introduction}
% \begin{verbatim}
%     Use command \yash{}, \sahana{}, \ambuj{} for comments
% \end{verbatim}

Natural phenomena and engineered systems are often governed by ordinary and partial differential equations. Solving these equations enables a variety of tasks - predicting the evolution of the system through simulation (e.x., forecasting weather using Navier Stokes Equation) \citep{kalnay2003atmospheric,bauer2015quiet,staniforth1991semi}, addressing control and optimization problems (e.x., optimizing heat shield design for spacecrafts using heat transfer equations) \citep{anderson1989hypersonic,troltzsch2010optimal}, and tackling inverse problems based on external measurements (e.x. reconstructing brain activity from EEG data using electrophysiological models) \citep{baillet2001electromagnetic,grech2008review}. 

Traditionally, PDEs are solved using finite difference methods \citep{smith1985numerical,leveque2007finite}, which discretize the spatial and/or temporal domain and approximate the solution by solving a system of equations at the discrete grid points. Similarly, finite element methods \citep{hughes2003finite,bathe2006finite,brenner2008mathematical} construct a system of equations based on fitted curves for each finite element and then rely on numerical linear algebra techniques--such as the Jacobi and Gauss-Seidel method \citep{saad2003iterative}--to compute the solution. Such iterative methods, however, suffer from high computational costs, especially when finely resolved grids are required--regardless of the sparsity of the coefficient matrix. Moreover, these solvers lack generalization across different initial conditions, boundary conditions, or forcing functions, as even minor changes to any of these parameters require re-solving the entire system of equations from scratch. 

These challenges have motivated the development of neural operators \citep{kovachki2023neural} - a class of machine learning models that aim to learn the solution operator directly, enabling fast inference across a range of parameters. Neural operators lift the classical linear layer in neural networks to a linear operator--typically a kernel integral operator--between infinite-dimensional function spaces. While the universal approximation theorem of operators \citep{chen1995universal} suggests that neural operators can approximate solution operators with high accuracy, they are prone to spectral bias \citep{liu2021multiscale,liu2024mitigating,you2024mscalefno,khodakarami2025mitigating,xu2025understanding}--similar to traditional neural networks \citep{rahaman2019spectral,xu2019frequency}--and favor learning the low-frequency components of the solution operator, often struggling to capture high-frequency features that iterative solvers excel at. 

Recognizing this complementarity, HINTS \citep{zhang2024blending} proposes a hybrid solver that leverages the advantages of traditional iterative methods and neural operators  to learn a wide range of eigenmodes more uniformly and quickly. It achieves this by periodically incorporating predictions from a neural operator into iterative updates of the traditional solver. While HINTS observes significant empirical performance improvements in convergence and accuracy when compared to using the standalone neural operator and Jacobi solver, the algorithm's fixed number of solver iterations between neural operator passes can be viewed as wasteful if the high frequency errors are overly dampened, or insufficient if not dampened enough. Furthermore, when HINTS employs multigrid methods, they are often restricted to cycle patterns, i.e. V cycle or W cycle, of fixed depths making it difficult to adapt cycle patterns in response to the distribution of errors across the learned eigenmodes.
% We propose a model-based router that

We present a Greedy Hybrid Solver that, at each iteration, myopically selects a solver from a suite of available methods to reduce the error between the true solution and prediction. Assuming that the true error is known at every iteration, we show that a greedy approach to selecting a sequence of solvers, aimed at minimizing prediction error, achieves at least a $1- e^{-\alpha}$-approximation to the optimal strategy. Given that true error values are typically unavailable at inference time, a convex surrogate loss is introduced using principles of learning to defer,  which, when minimized, enables the learned model to imitate the greedy solver without access to true error information. This alignment is guaranteed through the property of Bayes consistency. We empirically demonstrate that our approximate greedy solver outperforms HINTS and individual solvers in terms of uniform convergence across frequencies, on benchmark PDEs including the Poisson Equation, Helmholtz Equation, and the diffusion equation.


\section{Problem Setting and Background} \label{sec:background}

Let $D \subseteq \mathbb{R}^n$ denote the spatial domain, and let $\mathcal{A} = \{a: D \to \mathbb{R}^d\}$, $\mathcal{F} = \{f: D \to \mathbb{R}^o\}$, and $\mathcal{U} = \{u: D \to \mathbb{R}^m\}$ be the Banach space of coefficient, forcing, and solution functions, respectively.
% Let $\mathcal{A} = \{a: D \to \mathbb{R}^d\}$ be an infinite-dimensional Banach space of coefficient functions, $\mathcal{F} = \{f: D \to \mathbb{R}^o\}$ be an infinite-dimensional Banach space of forcing functions and $\mathcal{U} = \{u: D \to \mathbb{R}^m\}$ be the infinite dimensional Banach space of output functions where $D \subseteq \mathbb{R}^n$ is the spatial domain.
Let's consider the following Linear PDE:
\begin{equation} \label{eq:mainpde}
    \begin{split}
    \mathcal{L}_{x}^a\left(u\right) &= f ,x \in D \\
    \mathcal{B}_{x}\left(u\right) &= g ,x \in \partial D
    \end{split}
\end{equation}
% $$\mathcal{L}_{x}\left(u;k\right) = f$$
where $\mathcal{L}_x^a$ is a differential operator with respect to the spatial variable $x \in D$ parameterized by $a \in \mathcal{A}$,  $f \in \mathcal{F}$ is an external forcing function,  $u \in \mathcal{U}$ is the solution to the PDE, $\mathcal{B}_{x}$ is the boundary operator, and $g$ is the boundary term .  

PDEs of this form can be solved using a variety of numerical methods, including finite difference, finite element, spectral methods \citep{boyd2001chebyshev,canuto2006spectral}, and neural operators. In this paper, however, we focus specifically on establishing efficient collaborations between iterative or multigrid solvers from finite difference discretizations and neural operators.

\subsection{Finite Difference Methods} \label{sec:finitedifference}

Finite difference methods \citep{smith1985numerical,leveque2007finite} represent a class of numerical methods that approximate solutions to partial differential equations by replacing partial derivatives with discrete differential quotients. For example, the first derivative can be approximated as $\frac{\partial u(x)}{\partial x} \approx \frac{u(x + h) - u(x)}{h}$. When the entire domain is discretized,  evaluating the solution to the PDE $u$ at these grid points reduces to solving a system of algebraic equations or $Au = b$. While direct methods like Gauss Elimination and Thomas Algorithm are often employed to approximate the solution to such systems of equations, they can be computationally expensive in high-dimensional domains. In contrast, iterative methods like Jacobi and Gauss-Seidel provide computational speedups by iteratively updating the solution, gradually converging to the true solution for the PDE. The general iterative update can be written as:
\begin{equation}\label{eq:iterativeupdate}
    u^{(t + 1)} = u^{(t)} + C\left(b - Au^{(t)}\right)
\end{equation}
where $u^{(t + 1)}$ is the $(t + 1)$ iterate of the solution and $C$ is a preconditioning matrix. For the Jacobi method, $C_\mathrm{Jac} = D^{-1}$, where $D$ is the diagonal of $A$; for Gauss-Seidel, $C_\mathrm{GS} = (D + L)^{-1}$, where $L$ denotes the strict lower triangular part. 

However, convergence of these basic iterative methods is often slower for the low frequency components of the solution. Multigrid solvers \citep{briggs2000multigrid,trottenberg2001multigrid,hackbusch2013multi} are an advanced class of iterative methods that tackle this problem with specialized preconditioning matrices  Specifically, the solvers allow high frequency error components to be sufficiently dampened on a finer discretization and leaving the smooth parts of the error to a coarser grid where low frequency function appear high frequency. We defer the reader to \Cref{sec:multigrid} for more background on Multigrid solvers and the particular form of the preconditioning matrix.

\sahana{Might move the next two paragraphs to the appendix}

Let $A_h, A_{2h}$ represent the coefficient matrix on a fine grid with discretization parameter $h$ and $2h$, $u_h$ and $f_h$ represent the PDE solution and constant vector on a grid discretized by $h$, $R^{2h}_h$ denote the restriction matrix which transfers vectors from a fine grid to a coarse one, and $I^{h}_{2h}$ is the interpolation matrix transfers vectors from a coarse grid to a fine one. In a 2-grid method, a few iterations of the smoother (e.x. Jacobi or Gauss-Seidel) are first applied on the fine grid to approximate the solution of $A_hu_h = f_h$. The residual is then computed as $r_h = f_h - A_hu_h$ and restricted to the coarse grid via $r_{2h} = R^{2h}_hr_h$. The error equation, $A_{2h}e_{2h} = r_{2h}$, is solved on the coarse grid. The resulting estimate of the error is interpolated in the fine grid by $e_h = I^{h}_{2h}e_{2h}$, and the fine grid solution is updated by adding this correction, $u_h = u_h + e_h$. Finally, additional smoothing steps are performed on the fine grid to further reduce any high frequency errors. The preconditioning matrix for two-grid solver is $C_\mathrm{2G} = I^{h}_{2h}A_{2h}^{-1}R^{2h}_h$.

More complex strategies for multigrid like V-cycle and W-cycle compute error corrections recursively across multiple grids of varying coarseness (See \Cref{sec:multigrid} for pseudocode of these methods).

\subsection{Neural Operators} \label{sec:neuraloperators}

Suppose we have some observations $\{(a_i, f_i, u_i)\}_{i = 1}^N$ such that $(a_i, f_i) \sim \mathcal{P}_{\mathcal{A} \times \mathcal{F}}$ are i.i.d. samples drawn from some distribution. Let $u_i$ is generated using a deterministic solution operator called $\mathcal{G}^{*} : \mathcal{A} \times \mathcal{F} \to \mathcal{U}$ such that $u_i = \mathcal{G}^{*}(a_i, f_i)$. Assuming a Lebesgue measure over the domain $D$, for a given pair of $(a_i, u_i, f_i)$, neural operators  $\mathcal{G}_{\theta}$ seek to approximate the true solution operator $\mathcal{G}^*$ by considering the following loss:
\begin{equation} \label{eq:NOloss}
    l_{\text{NO}}(a_i, f_i, u_i, \mathcal{G}_{\theta}) = \left\|u_i - \mathcal{G}_{\theta}(a_i, f_i)\right\|^2_{\mathcal{U}} = \int_{D} \left\|u_i(x) - \mathcal{G}_{\theta}(a_i, f_i)(x)\right\|_2^2dx
\end{equation}

The risk of this loss function is denoted by $\mathcal{R}_{\text{NO}}(\mathcal{G}_{\theta}) = \mathbb{E}_{(a, f) \sim  \mathcal{P}_{\mathcal{A} \times \mathcal{F}}}\left[l_{\text{NO}}(a, f, \mathcal{G}^{*}(a, f), \mathcal{G}_{\theta})\right] = \mathbb{E}_{(a, f) \sim  \mathcal{P}_{\mathcal{A} \times \mathcal{U}}}\left[l_{\text{NO}}(a, f, u, \mathcal{G}_{\theta})\right]$. Ideally, $\mathcal{G}_{\theta}$ would minimize $\mathcal{R}_{\text{NO}}(\mathcal{G}_{\theta})$.

In classical PDE theory, if $\mathcal{L}_{x}^a$ is a linear operator, the Green's function \citep{stakgold2011greens} $G: D \times D \to \mathbb{R}^m$ is the solution for $\mathcal{L}_{x}^aG(x, .) = \delta_x$ and the solution function takes the following form: $u(x)_j =\int_D G(x, y)_j f(y)_jdy$. Inspired by this, neural operators \citep{kovachki2023neural} replace the traditional linear layer with a linear operator in the form of a kernel integral transform, $(\mathcal{K}v)(x) = \int_{D}\kappa(x, y)v(y)d\nu_t(y)$, mimicking the actions of Green's functions in a learnable way.  A single neural operator layer that maps a function $v_t: D \to \mathbb{R}^{d_t}$ to function $v_{t + 1}: D \to \mathbb{R}^{d_{t + 1}}$. This mapping takes the following form: $$v_{t + 1}(x) = \sigma_t\left(W_tv_t(x) + (\mathcal{K}_tv_t)(x) + b_t(x)\right)$$ for every $x \in D_{t + 1}$ where $W_t \in \mathbb{R}^{d_{t + 1} \times d_{t}}$ are matrices performing a local linear operation (discretization invariant), $\mathcal{K}_t: \{v_t: D \to \mathbb{R}^{d_{t + 1}}\} \to \{v_{t + 1}: D \to \mathbb{R}^{d_{t + 1}}\}$ are global kernel integral operators, $b_t: D \to \mathbb{R}^{d_{t + 1}}$ are bias functions which are often constant, and $\sigma_t$ is a fixed pointwise activation function mapping from $\mathbb{R}^{d_{t + 1}}$ to $\mathbb{R}^{d_{t + 1}}$.


The architecture of the neural operator is discretization-invariant, a property that is clear in the parameterization of $W_t$ and $b_t$. However, a variety of discretization-invariant kernel parameterizations and kernel integral approximation strategies have given rise to many different variants of neural operators.  
% Different variants of neural operators arise from the choice of kernel parameterization and the approximation strategy used to compute the integral. 
For example, Fourier Neural Operators \citep{li2020fourier} parameterize the kernel in the Fourier domain and use convolutions to characterize the integral. Graph Neural Operators \citep{li2020multipole} compute the integral over a subdomain using the Nystrom's approximation. Prior to the rise of neural operators, DeepONets \citep{lu2021learning} were a popular deep learning approach to operator learning. They can also be viewed a neural operator that uses a pointwise kernel with a discretized integral operator. 




\subsection{Spectral Bias} \sahana{might remove or move to the appendix}

Spectral bias \citep{rahaman2019spectral}, also known as F-principle \citep{xu2019frequency}, refers to the tendency of neural networks to learn low-frequency components of a function faster than high-frequency components during training. While the original authors for the spectral bias and F-principle primarily studied empirical manifestations of this phenomenon, theoretical work \citep{luo2019theory} shows that networks with arbitrary number of layers and width initially have larger loss gradients for low-frequency components with respect to the Fourier basis, causing them to prioritize learning the smooth components of the true mapping. 

Even in the Neural Tangent Kernel regime where networks are infinitely wide, spectral bias has been proven to persist \citep{cao2019towards}. This has also been observed in neural operators \citep{liu2021multiscale,liu2024mitigating,you2024mscalefno,khodakarami2025mitigating,xu2025understanding} where during early training, the model learns to output low-frequency parts of the output functions. A spectral decomposition of the loss described in \Cref{eq:NOloss} makes it easy to disentangle the errors associated with low- and high-frequency parts. Regardless of the exact $a$ and $f$, the solution space of this PDE takes on general form - $u(x) = \sum_{k}\widehat{u}(k) \phi_k(x)$ where $\{\phi_k(x)\}$ are the eigenfunctions of the Fourier basis and $\widehat{u}(k) = \left<u(x), \phi_k(x)\right> = \int_D u(x)\phi_k(x)dx$. These coefficients $\widehat{u}(k)$ vary with the different instantiations of $a$ and $f$. 

Accordingly, $l_{\text{NO}}$ can be expressed in terms of the eigenfunction coefficients. Using the orthonormality of the eigenfunctions, we obtain:
\begin{align*}
    l_{\text{NO}}(a_i, f_i, u_i, \mathcal{G}_{\theta}) &= \int_{D} \left\|\sum_{k = 1}^{\infty}\widehat{u}(k) \phi_k(x) - \sum_{k = 1}^{\infty}\widehat{u}_{\theta}(k) \phi_k(x)\right\|_2^2dx \\
    % &= \int_{D} \left\|\sum_{i = 1}^{\infty} \left(c_i -\widehat{c}_i \right)\phi_i(x)\right\|^2_2dx \\
    &= \int_{D} \sum_{k = 1}^{\infty} \left(\widehat{u}(k) -\widehat{u}_{\theta}(k) \right)^2\|\phi_k(x)\|^2_2 + \sum_{j \neq k} \left(\widehat{u}(j) -\widehat{u}_{\theta}(j) \right)\left(\widehat{u}(j) -\widehat{u}_{\theta}(j) \right)\phi_k(x)^{\top}\phi_j(x)dx \\
    % &=  \sum_{i = 1}^{\infty} \left(c_i -\widehat{c}_i \right)^2 \int_{D}\|\phi_i(x)\|^2_2dx + \sum_{j \neq i} \left(c_i -\widehat{c}_i \right)\left(c_j -\widehat{c}_j \right)^2 \int_{D}\phi_i(x)^{\top}\phi_j(x)dx \\
    &= \sum_{k = 1}^{\infty} \left(\widehat{u}(k) -\widehat{u}_{\theta}(k) \right)^2
\end{align*}
\sahana{We can include a background section on greedy set optimization or consistency instead??}

\section{General framework for Hybrid Solvers}

\sahana{Feel like I should include this in \Cref{sec:background} instead}

Let $h$ be the grid size and let $G_h(D)$ be an uniform grid with spacing $h$ over the bounded domain $D \subseteq \mathbb{R}^n$. Given a function $g: D \to \mathbb{R}$, we represent its restriction to $G_h(D)$ by $g_h$, i.e. vector of $g$ evaluated at the points of $G_h(D)$.
% Let $h$ be a fixed grid size and let an $h$-grid over the bounded domain $D \subseteq \mathbb{R}^n$  be $G_h(D) := \{\mathbf{k}h : \mathbf{k} \in \mathbb{Z}^n, \mathbf{k}h \in D\}$. We denote a restriction of the function $g: D \to \mathbb{R}$ to the grid $G_h(D)$ as $g_h := g(x) \Big\vert_{G_h(D)}$ and this can be interpreted as $g_h := \begin{bmatrix}
%     g(x_1) & \cdots & g(x_{|G_h(D)|})
% \end{bmatrix}^{\top} \in \mathbb{R}^{|G_h(D)|}$ for $x_1, x_2, \cdots, x_{|G_h(D)|} \in G_h(D)$. 
If $N = |G_h(D)|$, the discretized version of \Cref{eq:mainpde} is expressed as:
\begin{equation} \label{eq:mainpdediscrete}
    \mathcal{L}_h^au_h = f_h
\end{equation}
where $\mathcal{L}_h^a$ is the discretized operator represented as an $N \times N$ matrix, reducing the PDE reduced to a linear system of equations. To solve \Cref{eq:mainpdediscrete}, consider the following general hybrid solver: at each time step $t + 1$, the update is given by
\begin{equation}
     u^{\left(t + 1\right)}_h = u^{\left(t\right)}_h + C_{S_t}\left(f_h - \mathcal{L}_h^au^{\left(t \right)}_h\right)
\end{equation}
where $\mathcal{C} =\{C_{j}\}_{j = 1}^K$ denotes the set of $K$ possible preconditioning matrices and $S_t \in [K] = \{1, \dots, K\}$. Each $C \in \mathcal{C}$ corresponds to a particular solver, such as Jacobi, Gauss-Seidel, Multigrid, or Neural Operators. This update naturally generalizes the form of classical iterative methods, as seen in \Cref{eq:iterativeupdate}, with the distinction that the preconditioning matrix can be adaptively chosen at each iteration. Furthermore, HINTS \citep{zhang2024blending} can be viewed as a special case of this hybrid solver where the hybrid solver has access to two solvers ($K = 2$):  $C_1$ is the preconditioner for the neural operator, and $C_2$ is a preconditioner for a traditional solver. The routing decision is determined by $S_t = \1_{t \bmod \tau > 0} + 1$ which selects $C_1$ every $\tau$ steps and $C_2$ otherwise.

% indexed by $j \in [K] = \{1, \dots, K\}$, and the index $j$ may change with time. 
% In the case of HINTS , t

If the error at time step $t + 1$ is denoted by $e^{(t + 1)}_h = u_h - u^{\left(t + 1\right)}_h$, the update rule can be written in terms of the error from the previous step:
\begin{align*}
    e^{(t + 1)}_h = u_h -  u^{\left(t\right)}_h -C_{S_t}\left(\mathcal{L}_h^au_h - \mathcal{L}_h^au^{\left(t\right)}_h\right) = \left(I_N - C_{S_t}\mathcal{L}_h^a\right)e^{(t)}_h 
    % &= e^{(t)} - C_{S_t}\left(\mathcal{L}_h^a\left(u_h-u^{\left(t\right)}_h\right)\right) \\
    %&= e^{(t)} - C_{S_t}\mathcal{L}_h^ae^{(t)}
\end{align*}
where $I_N$ is an $N \times N$ identity matrix and $\left(I_N - C_{j}\mathcal{L}_h^a\right)$ is the error propagation matrix for the $j^{\text{th}}$ solver.

% The error after $T$ steps becomes:
% \begin{align*}
%     e^{(T)}_h = \left(I_N - C_{S_T} \mathcal{L}_h^a\right) \cdots \left(I_N - C_{S_1} \mathcal{L}_h^a\right)  e^{(0)}_h 
% \end{align*}
% where $I_N$ is the identity matrix. 

The objective of the hybrid solver is to select a sequence of solvers at each time that minimize the error norm after $T$ steps or $\left\|e^{(T)}_h\right\|_2^2$. Formally, we seek to solve 
\begin{equation} \label{eq:errorobjective}
    \min_{S_t \in [K], |S| \leq T} h(S)
\end{equation}
where $S = (S_1, \dots, S_T)$ is a sequence of solver indices and the objective function $h(S)$ is defined as
\begin{equation} \label{eq:errorobjective2}
    h(S) = \left\|\prod_{t = |S|}^1\left(I_N - C_{S_t} \mathcal{L}_h^a\right)  e^{(0)}_h \right\|^2_2
\end{equation}
Here, the product is ordered from right to left (i.e. from $t = 1$ to $t = |S|$), reflecting the sequential application of the preconditioners. 

% If running each solver $C_{j}$ incurs a time cost $c_j$, we can formulate a budget constrained optimization problem
% \begin{align*}
%     \min_{S \in [1, K]^*} & h(S) \\
%     \text{subject to} \quad &\sum_{t = 1}^T c_{i_t} \leq B
% \end{align*}

% where $B$ is the total computational budget.

\section{Greedy Algorithm} \label{sec:greedy}

\Cref{eq:errorobjective} defines a non-convex and combinatorial optimization problem whose search space grows exponentially in the time horizon $T$, rendering exact solutions computationally intractable for large $T$ or $K$. However, if the objective function satisfies supermodularity and monotonicity, a greedy algorithm -- such as the one described in \Cref{alg:greedy} -- can achieve solutions provably close to the optimal.

\sahana{this Para below feels like a backgroundy paragraph}

Supermodularity, also known as the increasing differences property, characterizes the idea that appending an element earlier in a sequence has a larger marginal effect than doing so later. The suboptimality of the greedy algorithm under various structural conditions (i.e. submodularity, monotonicity) has been extensively studied for both minimization \citep{liberty2017greedy,bounia2023approximating}  and maximization \citep{nemhauser1978analysis,das2018approximate,feige2011maximizing,bian2017guarantees,harshaw2019submodular} of set functions. However, these classical results are limited to permutation-invariant objective functions, where the function value doesn't depend on the order of the elements. In contrast, results on the suboptimality of greedy sequence maximization \citep{streeter2008online,alaei2021maximizing,zhang2015string,bernardini2020through,van2024performance,tschiatschek2017selecting} -- where the ordering of the elements affects the function -- have led to various definitions of submodularity and monotonicity in sequential settings. In this work, we leverage some of those tools to analyze the suboptimality of greedy solution for sequence minimization such as \Cref{eq:errorobjective}.

\begin{algorithm} 
\caption{Greedy Algorithm}\label{alg:greedy}
\begin{algorithmic}
\Require $\{C_{j}\}_{j = 1}^{K}, T, \mathcal{L}_h^a, e^{(0)}_h$
\State $S^{(0)} \gets \emptyset$
\For{$t \leq T$}
    \State $S^{(t + 1)} \gets S^{(t)} \ \oplus \argmin_{j \in [K]}\left\|\left( I_N - C_{j} \mathcal{L}_{h}\right)e^{(t)}_h\right\|^2_2$
    \State $e^{(t + 1)}_h \gets \left(I_N - C_{S_{t+ 1}} \mathcal{L}_{h}^a\right)e^{(t)}_h$
\EndFor
\State \Return $S^{(T)}$
\end{algorithmic}
\end{algorithm}

\subsection{Suboptimality of Greedy Sequence Minimization}

Let $\Omega^*$ represent the space of sequences with each element in $\Omega$. Let $S = (S_1, \dots, S_n) \in \Omega^*$ and $S' = (S_1', \dots, S_m') \in \Omega^*$ be two sequences. We denote their concatenation as $S \oplus S' = (S_1, \dots, S_n, S_1', \dots, S_m')$. If $S' = S \ \oplus L$ where $L \in \Omega^*$, then $S$ is a prefix of $S'$ which is denoted as $S \preceq S'$. A sequence function $g: \Omega^* \to \mathbb{R}$ is considered prefix monotonically non-increasing if $g(S \ \oplus S') \leq g(S)$ for all $S, S' \in \Omega^*$. Similarly, $g$ is postfix monotonically non-increasing if $g(S' \ \oplus S) \leq g(S)$ for all $S, S' \in \Omega^*$. \sahana{I may not need this. will discuss later}$g$ is considered weak-$\mu$-postfix monotonic if for all $S, S' \in \Omega^*$, 
\begin{equation*}
    g(S' \ \oplus S) \leq \mu g(S)
\end{equation*}
$\mu$ is generally greater than $1$ and $\mu = 1$ for postfix monotonic functions.

A prefix non-increasing function $g$ is considered sequence supermodular if, for all $S', S \in \Omega^*$ such that $S \preceq S'$, it holds that
\begin{align*}
    g(S) - g(S \ \oplus \omega ) \geq g(S') - g(S' \ \oplus \omega )
\end{align*}
for any $\omega \in \Omega$. However, the objective function $h(S)$ described in \Cref{eq:errorobjective2} may not, in general, satisfy this property. Therefore, we introduce weakly-$\alpha$-supermodularity for sequences. A prefix non-increasing function $g$ is weakly-$\alpha$-supermodular if, for any $S, S' \in \Omega^T$ such that $|S'| = T$, 
\begin{equation*}
     g(S) - g(S\ \oplus S') \leq \alpha T \max_{i \in [|S'|]} g(S) - g(S \ \oplus S'_i)
\end{equation*}
where $\alpha \geq 1$. When $\alpha = 1$, the functions is sequence supermodular (See \Cref{sec:proofalpha=1} for proof)



\Cref{th:suboptgreedy} characterizes the suboptimality of greedy solutions to weakly-$\alpha$-supermodular, prefix monotonic, and weakly-$\mu$-postfix montonic sequence functions. 


\begin{theorem} \label{th:suboptgreedy}
    Let $g: \Omega^* \to \mathbb{R}$ be a prefix and $\mu$-postfix monotone sequence function with weak-$\alpha$-supermodularity. Let the greedy solution of length $T$ be $S^T$ and the optimal solution be $O$ such that $|O| \leq T$. Then,
    \begin{equation*}
        g(S^T) \leq  \mu\left(1 - \left(1 - \frac{1}{\alpha T}\right)^T\right)
        g(O) +  \left(1 - \frac{1}{\alpha T}\right)^Tf(\emptyset)
    \end{equation*}
    Equivalently, \sahana{stuff below is wrong}
    \begin{equation*}
        \frac{g(\emptyset) - g(S^T)}{g(\emptyset) - g(O)} \geq \mu\left(1 - \left(1 - \frac{1}{\alpha T}\right)^T\right) \geq \mu\left(1 - \frac{1}{e^{\alpha}}\right)
    \end{equation*}
\end{theorem}

The proof of the above theorem can be found in \Cref{sec:proofsuboptgreedy}. As $T \to \infty$, the factor $\left(1 - \left(1 - \frac{1}{\alpha T}\right)^T\right)$ in \Cref{th:suboptgreedy} decreases to $1 - e^{-\alpha}$, showing that the worst-case performance of the greedy algorithm saturates and does not degrade further with longer sequences. 

The applicability of \Cref{th:suboptgreedy} to designing a greedy hybrid solver depends on verifying that the sequence function $f$ satisfies both prefix, postfix monotonicity, and weak-$\alpha$-supermodularity-- properties established in Proposition \ref{th:weaklyalphasupermodular}. We defer the full proof of Proposition \ref{th:weaklyalphasupermodular} to \Cref{sec:proofweaklyalphasupermodular}.

\begin{proposition} \label{th:weaklyalphasupermodular}
    If $\|I_N - C_j\mathcal{L}_h^a\| \leq 1$ for all $j \in [K]$, $h$ is (i) prefix monotonically non-increasing and (ii) f is weakly-$\alpha$-supermodular with $$\alpha = \max_{S: |S| = T}\frac{\left\|I_N - \prod_{j = T}^{1}\left(I_N - C_{S_i}\mathcal{L}_h^a\right)\right\|}{1 - \max_{i \in [T]}\|I_N - C_{S_i}\mathcal{L}_h^a\| }$$Furthermore, if $I_N - C_j\mathcal{L}_h^a$ is invertible for all $j \in [K]$, $h$ is also postfix monotonically non-increasing.
\end{proposition}

Proposition \ref{th:weaklyalphasupermodular} thus ensures that all the conditions required by \Cref{th:suboptgreedy} are satisfied as long as no solver in $\mathcal{C}$ amplifies errors. Notably, any well-trained neural operator is expected to reduce the error norm, and traditional solvers such as Jacobi and Gauss-Seidel are known to exhibit this error-reducing property for linear elliptic PDEs under standard conditions. 

\Cref{th:suboptgreedy} also indicates that sequences exhibiting near-supermodular behavior, i.e. $\alpha \approx 1$ achieves the best lower bound. A natural case of this occurs when the PDE of interest is linear with constant coefficients and periodic boundary conditions. Under these conditions, the error propagation matrices of classical solvers such as Jacobi and Gauss-Seidel exhibit a circulant structure making them both diagonalizable by the discrete Fourier transform matrix. Furthermore, when the learned neural operator is a single-layer Fourier Neural Operator without non-linearities, it too is a diagonalizable by the discrete Fourier basis. More generally, when the the error propagation matrices $\left(I - C_j \mathcal{L}_{h}\right)$ are simultaneously diagonalizable, i.e each admits a spectral decomposition with a shared eigenbasis, the original problem reduces to an allocation problem. This reduction is critical in analyzing the supermodularity of the objective with simultaneously diagonalizable error propagation matrices.
\begin{proposition} \label{th:simultaneoussupermodular}
    Let $\|I_N - C_j\mathcal{L}_h^a\| \leq 1$ for all $j \in [K]$ and $\left(I - C_j \mathcal{L}_{h}^a\right) = P\Lambda_jP^{-1}$. Then, $h$ is supermodular. 
\end{proposition}
The proof of \Cref{th:simultaneoussupermodular} can be found in \Cref{sec:proofsimultaneoussupermodular}.



\begin{corollary}[Suboptimality of HINTS]
    
\end{corollary}


\section{Approximate Greedy Router} \label{sec:approxgreedy}

The results from \Cref{sec:greedy} indicate the error reduction achieved with the  greedy solution closely matches the of the optimal sequence. However, this greedy algorithm is predicated on having an accurate estimate of the initial error or $e^{(0)}$. A poor approximation of $e^{(0)}$ can result in miscalibrated decisions, with the approximation errors propagating and amplifying over subsequent steps. 

To remedy this, we wish to design a router that learns to select solvers myopically, as described in \Cref{alg:greedy}, without access to the true initial error $e^{(0)}$. At each time step $t$, the router $r$ selects a solver based on the coefficient function $a \in \mathcal{A}$, forcing function $f \in \mathcal{F}$ and the current iterate of the solution function $u^{(t)} \in \mathcal{U}$. Assuming that the initial guess $u^{(0)}$ is a deterministic function of $a$ and $f$, it follows that the entire sequence $\{u^{(t)}\}$ is fully determined by $a$ and $f$, even though this dependence is implicit.
% $u^{(t)}$ is a function of the initial guess of the solution function $u^{(0)} \sim \mathcal{P}_{\mathcal{U}}$ and 
When $r(a, f, u^{(t)}) = j$, the solution function is passed onto the $j^{\text{th}}$ solver and this results in an error of $\left\|\left( I_N - C_{j} \mathcal{L}_{h}\right)e^{(t)}\right\|^2_2$. Learning such a router requires minimizing the following loss objective

\begin{equation} \label{eq:routerloss}
    l_{\text{route}}\left(r,\mathcal{C}, a_h, f_h, u^{(t)}_h, u_h\right) = \sum_{j = 1}^K \left\|\left( I_N - C_{j} \mathcal{L}_{h}\right)\left(u_h - u^{(t)}_h\right)\right\|^2_2\1_{r(a_h, f_h, u^{(t)}_h) = j}
\end{equation}

The corresponding $l_{\text{route}}$-risk is defined as $\mathcal{R}_{\text{route}}\left(r,\mathcal{C}\right) = \mathbb{E}_{a, f \sim \mathcal{P}_{\mathcal{A} \times \mathcal{F}}}\left[l_{\text{route}}\left(r,\mathcal{C}, a_h, f_h, u^{(t)}_h, u_h\right)\right]$. Minimizing $\mathcal{R}_{\text{route}}\left(r\right)$ yields the following Bayes-Optimal Router:
\begin{align*}
    r^*(a_h, f_h, u^{(t)}_h) = \text{argmin}_{j \in [1, K]} \left\|\left( I_N - C_{j} \mathcal{L}_{h}^a\right)e^{(t)}_h\right\|^2_2
\end{align*}
\sahana{is the above equation clear or does this need to be proved somehow}
which aligns with the decision rule in \Cref{alg:greedy}, confirming that $l_{\text{route}}$ is consistent with learning the greedy algorithm.

\subsection{Surrogate Loss}

Since \Cref{eq:routerloss} is discontinuous and non-convex, direct minimization is intractable in practice. To overcome this, we introduce a surrogate loss that satisfies two key properties: (1) it is convex, enabling efficient optimization, and (2) it is Bayes consistent, ensuring the Bayes optimal decision of the original loss is preserved upon minimization. Formally, a surrogate $\phi$ is considered to be Bayes consistent with respect to the loss $l$ if 
\begin{equation*}
    \lim_{n \to \infty} \mathcal{R}_{\phi}(f_n) - \mathcal{R}_{\phi}^* \implies \lim_{n \to \infty} \mathcal{R}_{l}(f_n) - \mathcal{R}_{l}^* 
\end{equation*}
where $\mathcal{R}_{l}(f) = \mathbb{E}\left[l(f(X), Y)\right]$ and $\mathcal{R}_{l}^* = \inf_f\mathcal{R}_{l}(f)$. In other words, in the limit of infinite data, if a sequence of learned hypotheses $\{f_n\}$ converges to the optimal risk under $\phi$, it also converges to the optimal risk with respect to the original loss $l$.

To define a surrogate loss for the routing problem, consider a set of scoring functions $\mathbf{g} =\{g_j\}_{j = 1}^K$ with $g_j: \mathcal{A} \times \mathcal{F} \times \mathcal{U} \to \mathbb{R}$ and we define the router as $r(a, f, u^{(t)}) = \text{argmax}_{j \in [1, K]} g_j(a, f, u^{(t)})$. Then, minimizing the following surrogate loss yields the same decision as minimizing \Cref{eq:routerloss}:
\begin{equation} \label{eq:routersurrogate}
    \Psi\left(\mathbf{g}, \mathcal{C}, a, f, u^{(t)}, u\right) = -\sum_{j = 1}^K \left(\Tilde{c}\left(a, f, u^{(t)}, u\right) -\left\|\left( I_N - C_{j} \mathcal{L}_{h}\right)e^{(t)}\right\|^2_2\right) \log\left(\frac{e^{g_j(a, f, u^{(t)})}}{\sum_{k = 1}^Ke^{g_k(a, f, u^{(t)})}}\right)% \psi\left(\mathbf{g}(a, f, u^{(t)}),j\right)
\end{equation}

where 
% $\psi\left(\mathbf{g}(a, f, u^{(t)}),j\right)$ is a consistent surrogate for the multiclass $0/1$ loss and 
$\Tilde{c}\left(a,f, u^{(t)}, u\right) = \max_{j \in [1, K]} \left\|\left( I_N - C_{j} \mathcal{L}_{h}\right)e^{(t)}\right\|^2_2$. The $\Psi$-risk is denoted by $\mathcal{R}_{\Psi}\left(\mathbf{g}, \mathcal{C}\right) = \mathbb{E}_{a, f \sim \mathcal{P}_{\mathcal{A} \times \mathcal{F}}}\left[\Psi\left(\mathbf{g}, \mathcal{C}, a, f, u^{(t)}, u\right)\right]$. 

The convexity of $\Psi$ with respect to $\mathbf{g}$ follows from the fact that the log-softmax function is convex in its inputs. \Cref{th:consistency} shows that $\Psi$ achieves bayes consistency with respect to $l$

\begin{theorem} \label{th:consistency}
    Let $\left\|\left( I_N - C_{j} \mathcal{L}_{h}^a\right)e^{(t)}_h\right\|^2_2 < \Bar{E} < \infty$  for all $j \in [1, K]$. If there exists $j, k \in [1, K]$ such that $\left|\left\|\left( I_N - C_{k} \mathcal{L}_{h}^a\right)e^{(t)}_h\right\|^2_2 - \left\|\left( I_N - C_{j} \mathcal{L}_{h}^a\right)e^{(t)}_h\right\|^2_2\right| > \Delta > 0$, then, for any collection of solvers $\{C_j\}_{j= 1}^K$, $\Psi$ is Bayes consistent surrogate for $l_{\text{route}}$.
\end{theorem}

\sahana{maybe a generalization bound here}

\section{Related Works}

\sahana{not sure what to include here }

\section{Experiments}

\section{Discussion}


\section{Random Notes}

\begin{itemize}
    \item Difficulty of finite sample guarantees: If we know that $f(S \ \oplus r(S)) - \min_{\omega \in \Omega} f(S \ \oplus \omega) < \epsilon$ (can obtain this lower bound from the consistency bound), this doesn't guarantee that $f(Z^T)$ where $Z^T$ is constructed using $r$ such that $Z^T = r(\varnothing) \ \oplus r(Z_1) \ \oplus r(Z_2) \ \oplus \ \cdots \oplus r(Z_{T - 1})$ is lower bounded by the performance of the true greedy sequence $S^T$,
    \item Budget constraints: There is some work on how to 
    \item $f$ is not needed for the router
    \item can neural operator be a matrix?
    \item check if the error norms are not spctral radius
    \item TODO: include proof of consistency, generalization bounds (maybe)
    \item TODO: look up this upper bound weirdness for the surrogate
    \item \Cref{sec:approxgreedy} I am not sure if neural operators mathematically act on function spaces. I have the routers recieving discretized inputs mathmatically so not sure if it's consistent with neural operator literature. but we do need to connect \Cref{sec:greedy} somehow
    \item what is $\mu$
\end{itemize}

\bibliography{iclr2026_conference}
\bibliographystyle{iclr2026_conference}

\appendix
\section{Background: Multigrid} \label{sec:multigrid}

\section{Proofs for \Cref{sec:greedy}} \label{sec:proofgreedy}

\subsection{Proof of Proposition \ref{th:alpha=1}} \label{sec:proofalpha=1}

\begin{proposition} \label{th:alpha=1}
    Any prefix monotonically non-increasing sequence supermodular function $g$ is weakly-$\alpha$-supermodular with $\alpha = 1$
\end{proposition}

\begin{proof}

This proof is adapted from \citet{liberty2017greedy}.

If $g$ is sequence supermodular then,
\begin{align*}
    &g(S) - g(S \ \oplus S' ) \\
    &= \sum_{i = 1}^{|S'|} g(S \ \oplus \left(S'_1, \dots, S'_{i - 1}\right)) - g(S \ \oplus \left(S'_1, \dots, S'_{i}\right)) \\
    &\overset{(a)}{\leq} \sum_{i = 1}^{|S'|} g(S) - g(S \ \oplus  S'_{i}) \\
    &\leq |S'|\max_{i \in [|S'|]} g(S) - g(S \ \oplus  S'_{i})
\end{align*}

(a) by supermodularity


\end{proof}

\subsection{Proof of \Cref{th:suboptgreedy}} \label{sec:proofsuboptgreedy}



\begin{customtheorem}{\ref{th:suboptgreedy}}
    Let $g: \Omega^* \to \mathbb{R}$ be a prefix and $\mu$-postfix monotone sequence function with weak-$\alpha$-supermodularity. Let the greedy solution of length $T$ be $S^T$ and the optimal solution be $O$ such that $|O| \leq T$. Then,
    \begin{equation*}
        g(S^T) \leq  \mu\left(1 - \left(1 - \frac{1}{\alpha T}\right)^T\right)
        g(O) +  \left(1 - \frac{1}{\alpha T}\right)^Tf(\emptyset)
    \end{equation*}
    Equivalently, \sahana{stuff below is wrong}
    \begin{equation*}
        \frac{g(\emptyset) - g(S^T)}{g(\emptyset) - g(O)} \geq \mu\left(1 - \left(1 - \frac{1}{\alpha T}\right)^T\right) \geq \mu\left(1 - \frac{1}{e^{\alpha}}\right)
    \end{equation*}
\end{customtheorem}

\begin{proof}

This proof strategy is inspired by \citet{streeter2008online} and \citet{liberty2017greedy}

\begin{align*}
    g(S^t) -g(O) &\overset{(a)}{\leq} g(S^t) - \frac{1}{\mu}g(S^t\ \oplus O) \\
    &= \left(1 - \frac{1}{\mu}\right)g(S^t) + \frac{1}{\mu}\left(g(S^t) - g(S^t\ \oplus O)\right) \\
    &=\left(1 - \frac{1}{\mu}\right)g(S^t) + \frac{1}{\mu}\sum_{i = 1}^{|O|} g(S^t \oplus \left(o_1, \dots o_{i - 1}\right)) -g(S^t\ \oplus \left(o_1, \dots, o_{i}\right))\\
    &\overset{(b)}{\leq} \left(1 - \frac{1}{\mu}\right)g(S^t) + \frac{\alpha}{\mu}\sum_{i = 1}^{|O|} g(S^t) -g(S^t\ \oplus  o_{i}) \\
    &\leq \left(1 - \frac{1}{\mu}\right)g(S^t) + \frac{\alpha|O|}{\mu}\max_{i \in [|O|]} g(S^t) -g(S^t\ \oplus  o_{i}) \\
    &\leq  \left(1 - \frac{1 - \alpha T}{\mu}\right)g(S^t) -  \frac{\alpha T}{\mu}\min_{\omega \in \Omega} g(S^t\ \oplus  \omega) \\
    &=  \left(1 - \frac{1 - \alpha T}{\mu}\right)g(S^t) -  \frac{\alpha T}{\mu}g(S^{t + 1})
    % &= Tg(S^t) -  Tg(S^{t + 1}) \\
    % &= T \left(g(S^t) - g(O)\right) - T\left(g(S^{t + 1}) -g(O)\right) \\
    % g(S^{t + 1}) -g(O) &\leq \left(1 - \frac{1}{T}\right)\left(g(S^t) - g(O)\right)
\end{align*}
(a) by $\mu$- postfix monotonicity, (b) by supermodularity. 

After rearranging the inequality, we get:

\begin{align*}
    g(S^{t + 1}) &\leq \frac{\mu}{\alpha T}\left(g(O) -  \frac{1 - \alpha T}{\mu}g(S^t)\right) \\
    &= \frac{\mu}{\alpha T}g(O) + \left(1 - \frac{1}{\alpha T}\right)g(S^t)
\end{align*}

When recursively applying this inequality, we get:

\begin{align*}
     g(S^{T}) &\leq \frac{\mu}{\alpha T}g(O)\sum_{i = 0}^{T - 1}\left(1 - \frac{1}{\alpha T}\right)^i + \left(1 - \frac{1}{\alpha T}\right)^Tg(\emptyset) \\
     &= \mu \left(1 - \left(1 - \frac{1}{\alpha T}\right)^T\right)g(O) + \left(1 - \frac{1}{\alpha T}\right)^Tg(\emptyset)
\end{align*}
\end{proof}

\subsection{Proof of Proposition \ref{th:weaklyalphasupermodular}} \label{sec:proofweaklyalphasupermodular}

\begin{customproposition}{\ref{th:weaklyalphasupermodular}}
    If $\|I_N - C_j\mathcal{L}_h^a\| \leq 1$ for all $j \in [K]$, $h$ is (i) prefix monotonically non-increasing and (ii) f is weakly-$\alpha$-supermodular with $$\alpha = \max_{S \in [K]^*: |S| = T}\frac{\left\|I_N - \prod_{t = T}^{1}\left(I_N - C_{S_t}\mathcal{L}_h^a\right)\right\|}{T\left(1 - \max_{t \in [T]}\|I_N - C_{S_t}\mathcal{L}_h^a\|\right) }$$Furthermore, if $I_N - C_j\mathcal{L}_h^a$ is invertible for all $j \in [K]$, $h$ is also postfix monotonically non-increasing.
\end{customproposition}

\begin{proof}

\textbf{Prefix Montonicity: } Let $S \preceq S'$ where $S' = S \ \oplus N$.
\begin{align*}
    h(S') &= \left \|\prod_{t = |S'|}^1\left(I_N - C_{S'_t}\mathcal{L}_h^a\right) e^{(0)}\right\|^2_2 \\
    &= \left \|\prod_{t = |S'|}^{|S| + 1}\left(I_N - C_{S'_t}\mathcal{L}_h^a\right)\prod_{t = |S|}^1\left(I_N - C_{S'_t}\mathcal{L}_h^a\right) e^{(0)}\right\|^2_2 \\
    &\leq \prod_{t = |S'|}^{|S| + 1}\left \|\left(I_N - C_{S'_t}\mathcal{L}_h^a\right) \right\|^2_2 \left\|\prod_{t = |S|}^1\left(I_N - C_{S_t}\mathcal{L}_h^a\right) e^{(0)}\right\|^2_2 \\
    &\leq h(S)
\end{align*}
\textbf{Weakly-$\alpha$-supermodular: } 
% \paragraph{ } 
We will upper bound $\alpha$ by providing an upper bound for $ f(S) - f(S\ \oplus S') $ and a lower bound for $\max_i f(S) - f(S\ \oplus S'_i) $

\begin{align*}
    h(S) - h(S\ \oplus S') &=  \left \|\prod_{t = |S|}^1\left(I_N - C_{S_t}\mathcal{L}_h^a\right) e^{(0)}\right\|^2_2 -  \left \|\prod_{t = |S'|}^{1}\left(I_N - C_{S'_t}\mathcal{L}_h^a\right) \prod_{t = |S|}^1\left(I_N - C_{S_t}\mathcal{L}_h^a\right) e^{(0)}\right\|^2_2 \\
    &\leq \left \|\prod_{t = |S|}^1\left(I_N - C_{S_t}\mathcal{L}_h^a\right) e^{(0)}-\prod_{t = |S'|}^{1}\left(I_N - C_{S'_t}\mathcal{L}_h^a\right) \prod_{t = |S|}^1\left(I_N - C_{S_t}\mathcal{L}_h^a\right) e^{(0)}\right\|^2_2 \\
    &=  \left \| \left(I_N-\prod_{t = |S'|}^{1}\left(I_N - C_{S'_t}\mathcal{L}_h^a\right) \right) \prod_{t = |S|}^1\left(I_N - C_{S_t}\mathcal{L}_h^a\right) e^{(0)}\right\|^2_2 \\
    &\leq  \left \| I_N-\prod_{t = |S'|}^{1}\left(I_N - C_{S'_t}\mathcal{L}_h^a\right) \right\|^2_2\left\| \prod_{t = |S|}^1\left(I_N - C_{S_t}\mathcal{L}_h^a\right) e^{(0)}\right\|^2_2
\end{align*}

To lower bound $\max_i f(S) - f(S\ \oplus S'_i) $
\begin{align*}
    \max_i h(S) - h(S\ \oplus S'_i) &= \left \|\prod_{t = |S|}^1\left(I_N - C_{S_t}\mathcal{L}_h^a\right) e^{(0)}\right\|^2_2 - \max_i \left \|\left(I_N - C_{S_i'}\mathcal{L}_h^a\right)\prod_{t = |S|}^1\left(I_N - C_{S_t}\mathcal{L}_h^a\right) e^{(0)}\right\|^2_2  \\
    &\geq \left \|\prod_{t = |S|}^1\left(I_N - C_{S_t}\mathcal{L}_h^a\right) e^{(0)}\right\|^2_2 - \max_i \left \|I_N - C_{S_i'}\mathcal{L}_h^a\right\|\left\|\prod_{t = |S|}^1\left(I_N - C_{S_t}\mathcal{L}_h^a\right) e^{(0)}\right\|^2_2 \\
    &= \left(1 - \max_i \left \|I_N - C_{S_i'}\mathcal{L}_h^a\right\|^2_2\right) \left \|\prod_{t = |S|}^1\left(I_N - C_{S_t}\mathcal{L}_h^a\right) e^{(0)}\right\|^2_2 
\end{align*}

Finally,

\begin{align*}
    \frac{h(S) - h(S\ \oplus S')}{T\max_i h(S) - h(S\ \oplus S'_i)} &\leq \max_{S, S' \in \Omega^T, |S'| = T}\frac{ \left \| I_N-\prod_{t = |S'|}^{1}\left(I_N - C_{S'_t}\mathcal{L}_h\right) \right\|^2_2\left\| \prod_{t = |S|}^1\left(I_N - C_{S_t}\mathcal{L}_h\right) e^{(0)}\right\|^2_2}{T\left(1 - \max_i \left \|I_N - C_{S_i'}\mathcal{L}_h\right\|^2_2\right) \left \|\prod_{t = |S|}^1\left(I_N - C_{S_t}\mathcal{L}_h\right) e^{(0)}\right\|^2_2 } \\
    &= \max_{S' \in \Omega^T, |S'| = T}\frac{\left \| I_N-\prod_{t = |S'|}^{1}\left(I_N - C_{S'_t}\mathcal{L}_h\right) \right\|^2_2}{T - T\max_i \left \|I_N - C_{S_i'}\mathcal{L}_h\right\|^2_2}
\end{align*}

\textbf{Postfix Montonicity: } Let $S' = S \ \oplus N$.
\begin{align*}
    h(S') &= \left \|\prod_{t = |S'|}^1\left(I_N - C_{S'_t}\mathcal{L}_h\right) e^{(0)}\right\|^2_2 \\
    &= \left \|\prod_{t = |N|}^{1}\left(I_N - C_{N_t}\mathcal{L}_h\right)\prod_{t = |S|}^1\left(I_N - C_{S_t}\mathcal{L}_h\right) e^{(0)}\right\|^2_2 \\
    &\overset{(a)}{=}  \left \|\prod_{t = |N|}^{1}\left(I_N - C_{N_t}\mathcal{L}_h\right)\prod_{t = |S|}^1\left(I_N - C_{S_t}\mathcal{L}_h\right)\left(\prod_{t = |N|}^{1}\left(I_N - C_{N_t}\mathcal{L}_h\right)\right)^{-1}\prod_{t = |N|}^{1}\left(I_N - C_{N_t}\mathcal{L}_h\right) e^{(0)}\right\|^2_2 \\
    &\leq \left \|\prod_{t = |N|}^{1}\left(I_N - C_{N_t}\mathcal{L}_h\right)\right\|^2_2\left \|\prod_{t = |S|}^1\left(I_N - C_{S_t}\mathcal{L}_h\right)\right\|^2_2\left \|\prod_{t = |N|}^{1}\left(I_N - C_{N_t}\mathcal{L}_h\right)\right\|^{-2}_2\left \| \prod_{t = |N|}^{1}\left(I_N - C_{N_t}\mathcal{L}_h\right)e^{(0)}\right\|^2_2 \\
    &\leq \prod_{t = |S|}^{1}\left \|\left(I_N - C_{S_t}\mathcal{L}_h\right) \right\|^2_2 \left\|\prod_{t = |N|}^1\left(I_N - C_{N_t}\mathcal{L}_h\right) e^{(0)}\right\|^2_2 \\
    &\leq h(N)
\end{align*}

(a) due the invertibility of $I_N - C_j\mathcal{L}_h$

\end{proof}

\subsection{Proof of \Cref{th:simultaneousallocation}} \label{sec:proofsimultaneousallocation}

\begin{lemma} \label{th:simultaneousallocation}
     Let $\left(I - C_j \mathcal{L}_{h}^a\right) = P\Lambda_jP^{-1}$ where $P$ is an orthogonal matrix and $\Lambda_j = \mathrm{diag}(\lambda_{j1}, \dots, \lambda_{jN})$. If $P^{-1}e^{(0)}_h = z$, then the following equality holds:
     \begin{equation} \label{eq:simultaneous}
          h(S) = \sum_{i = 1}^N z_i^2 \prod_{j = 1}^K \lambda_{ji}^{2m_j(S)}
          % \qquad\mathrm{s.t}\qquad \sum_{k=1}^{K} m_k = T 
     \end{equation}
     where $m_j(S) = \sum_{t= 1}^{|S|}\1_{S_t = j}$
\end{lemma}

\begin{proof}

\begin{align*}
    h(S) &= \left\|\prod_{t = |S|}^1\left(I_N - C_{S_t} \mathcal{L}_{h}^a\right)  e^{(0)}_h \right\|^2 \\
    &=  \left\|\prod_{t = |S|}^1\left(P\Lambda_{S_t}P^{-1}\right)  e^{(0)}_h \right\|^2 \\
    &=  \left\|P\prod_{t = |S|}^1\Lambda_{S_t}P^{-1}  e^{(0)}_h \right\|^2 \\
    &= \left\|\prod_{t = |S|}^1\Lambda_{S_t}P^{-1}  e^{(0)}_h \right\|^2\\
    &=  \sum_{i = 1}^N z_i^2 \prod_{t = T}^1 \lambda_{S_ti}^2 \\
    &=  \sum_{i = 1}^N z_i^2 \prod_{j = 1}^K \lambda_{ji}^{2m_j(S)} 
    % &= \min_{\mathbf{m} \in\mathbb{N}^{K}}\quad \sum_{j = 1}^N z_i^2 \exp \left\{\sum_{k = 1}^K m_k \log \lambda_{kj}^{2}\right\}
    % \qquad\mathrm{s.t}\qquad \sum_{k=1}^{K} m_k = T
\end{align*}

\end{proof}

\subsection{Proof of \Cref{th:simultaneoussupermodular}} \label{sec:proofsimultaneoussupermodular}
\begin{customproposition}{\ref{th:simultaneoussupermodular}}
    Let $\|I_N - C_j\mathcal{L}_h^a\| \leq 1$ for all $j \in [K]$ and $\left(I - C_j \mathcal{L}_{h}^a\right) = P\Lambda_jP^{-1}$. Then, $h$ is supermodular. 
\end{customproposition}

\begin{proof}

Let $S \preceq S'$ where $S' = S \ \oplus B$. By \Cref{th:simultaneousallocation},
\begin{align*}
    h(S) &= \sum_{i = 1}^N z_i^2 \prod_{j = 1}^K \lambda_{ji}^{2m_j(S)} 
\end{align*}
where $m_j(S) = \sum_{t = 1}^{|S|} \1_{S_t = j}$ be the number of times a sequence $S$ calls the solver $j$. Recall that $h$ is considered sequence supermodular if $\forall S', S \in \Omega^*$ such that $S \preceq S'$, it holds that
\begin{align*}
    h(S) - h(S \ \oplus \omega ) \geq h(S') - h(S' \ \oplus \omega )
\end{align*}
% Without loss of generality, let's assume that $\omega$ that will be appended to $S, S'$ is $1$.
\begin{align*}
    h(S) - h(S \ \oplus \omega) &= \sum_{i = 1}^N z_i^2 \prod_{j = 1}^K \lambda_{ji}^{2m_j(S)} - \sum_{i = 1}^N z_i^2 \lambda_{\omega i}^{2}\prod_{k = 1}^K \lambda_{ji}^{2m_j(S)} \\
    &=  \sum_{i = 1}^N \left(1 - \lambda_{\omega i}^2\right) z_i^2 \prod_{j = 1}^K \lambda_{ji}^{2m_j(S)}
\end{align*}
Similarly, 
\begin{align*}
    h(S') - h(S \ \oplus \omega) &= \sum_{i = 1}^N \left(1 - \lambda_{\omega i}^2\right) z_i^2 \prod_{j = 1}^K \lambda_{ji}^{2m_j(S')} \\
    &\overset{(a)}{=}  \sum_{i = 1}^N \left(1 - \lambda_{\omega i}^2\right) z_i^2 \prod_{j = 1}^K \lambda_{ji}^{2\left(m_j(S) + m_j(B)\right)} \\
    &\overset{(b)}{\leq } \sum_{i = 1}^N \left(1 - \lambda_{\omega i}^2\right) z_i^2 \prod_{j = 1}^K \lambda_{ji}^{2m_j(S)} \\
    &= h(S) - h(S \ \oplus \omega)
\end{align*}


(a) Since $m_j(S') = \sum_{t = 1}^{|S'|} \1_{S_t' = j} = \sum_{t = 1}^{|S|} \1_{S_t' = j} + \sum_{t = |S|}^{|S'|} \1_{S_t' = j} = \sum_{t = 1}^{|S|} \1_{S_t = k} + \sum_{t = 1}^{|B|} 1_{B_t = j} = m_j(S) + m_j(B)$, (b) since $\rho(I_N - C_j \mathcal{L}_h^a) < 1$ and $m_j(B) \geq 0$ for all $j \in [K]$

\end{proof}


\section{Proofs for \Cref{sec:approxgreedy}}

\end{document}
